\subsection{INVARIANTS OF A FINITE LINEAR GROUP: RING-THEORETIC PROPERTIES}

THEOREM 3. We retain the assumptions and notation of Th. 1. In the set of systems of generators of the ideal \(R_{+}\) of R consisting of homogeneous elements, choose a minimal element \((\alpha_{1},\ldots,\alpha_{l})\) . Let \(k_{i}\) be the degree of \(\alpha_{i}\) . Assume that the \(k_{i}\) are coprime to the characteristic exponent of K. Then \(l=\dim V\) , the \(\alpha_{i}\) generate the K-algebra R, and are algebraically independent over K. In particular, R is a graded K-algebra of polynomials of transcendence degree l over K.

The assumption made about the \(k_{i}\) is superfluous, but is irrelevant for the applications to finite Coxeter groups, since then K = R. Cf. no. 5, where we shall give another proof of Th. 3.

Th. 3 follows from Prop. 2 (i), Th. 1 and the following lemma:

Lemma 1. Let K be a commutative field, S a graded K-algebra of polynomials, and R a graded subalgebra of S of finite type such that the R-module S admits a basis \((z_{\lambda})_{\lambda \in \Lambda}\) consisting of homogeneous elements. In the set of systems of generators of the ideal \(R_{+}\) of R consisting of homogeneous elements, choose a minimal element \((\alpha_{1}, \ldots, \alpha_{s})\) . Assume that, for all i, the degree \(k_{i}\) of \(\alpha_{i}\) is coprime to the characteristic exponent p of K. Then the \(\alpha_{i}\) generate the K-algebra R and are algebraically independent over K.

By Algebra, Chap. II, § 11, no. 4, Prop. 7, the assumption made about the \(\alpha_{i}\) is equivalent to saying that they are homogeneous and that their images in the K-vector space \(R_{+}/(R_{+})^{2}\) form a basis of this space. This condition is invariant under extension of the base field; we can thus reduce to the case where the latter is perfect.

The family \((\alpha_{1},\ldots,\alpha_{s})\) generates the algebra R by Commutative Algebra, Chap. III, §1, no. 2, Prop. 1. We argue by contradiction and assume that this family is not algebraically independent over K.

\begin{enumerate}
\def\labelenumi{\arabic{enumi})}
\tightlist
\item
  We show first of all that there exist families
\end{enumerate}

\[
\left(\beta_ {i}\right) _ {1 \leqslant i \leqslant s}, \quad \left(y _ {k}\right) _ {1 \leqslant k \leqslant r}, \quad \left(d _ {i k}\right) _ {1 \leqslant i \leqslant s, 1 \leqslant k \leqslant r}
\]

of homogeneous elements of S with the following properties:

\[
\beta_ {i} \in R \text {for all} i, \text {and the} \beta_ {i} \text {are not all zero};\tag{7}
\]

\[
\deg y _ {k} > 0 \text {   for   all   } k;\tag{8}
\]

\[
\alpha_ {i} = \sum_ {k = 1} ^ {r} d _ {i k} y _ {k} \quad \text { for   all } i;\tag{9}
\]

\[
\sum_ {i = 1} ^ {s} \beta_ {i} d _ {i k} = 0 \quad \text { for   all } k.\tag{10}
\]

Let \(X_{1},\ldots,X_{s}\) be indeterminates and give \(K[X_{1},\ldots,X_{s}]\) with the graded algebra structure for which \(X_{i}\) has degree \(k_{i}\) . There exist non-zero homogeneous elements \(H(X_{1},\ldots,X_{s})\) of \(K[X_{1},\ldots,X_{s}]\) such that \(H(\alpha_{1},\ldots,\alpha_{s})=0\) ; choose H to be of minimum degree. If \(\partial H/\partial X_{i}\neq0\) , the polynomial \(\frac{\partial H}{\partial X_{i}}(\alpha_{1},\ldots,\alpha_{s})\) is a non-zero homogeneous element of R; if \(p\neq1\) , H is not of the form \(H_{1}^{p}\) with \(H_{1}\in K[X_{1},\ldots,X_{s}]\) . Take

\[
\beta_ {i} = k _ {i} \frac {\partial \mathrm{H}}{\partial \mathrm{X} _ {i}} (\alpha_ {1}, \dots , \alpha_ {s}).
\]

Since K is perfect, the polynomials \(\partial H/\partial X_{i} \in K[X_{1},\ldots,X_{s}]\) are not all zero (Algebra, Chap. V, §1, no. 3, Prop. 4); in view of the assumption made about the \(k_{i}\) , neither are the \(\beta_{i}\) .

On the other hand, S can be identified with a graded algebra of polynomials \(K[x_{1},\ldots,x_{r}]\) for appropriate indeterminates \(x_{1},\ldots,x_{r}\) with suitable degrees \(m_{i}>0\) . Let \(D_{k}\) be the partial derivative with respect to \(x_{k}\) on S. Take \(d_{ik} = k_i^{-1}\mathrm{D}_k(\alpha_i)\) . Then the equality (10) holds because its left-hand side is \(\mathrm{D}_k(\mathrm{H}(\alpha_1,\dots ,\alpha_s))\) . On the other hand, if we put \(y_{1} = m_{1}x_{1},\ldots ,y_{r} = m_{r}x_{r}\) , the equality (9) follows from the equality (5) of no. 1.

\begin{enumerate}
\def\labelenumi{\arabic{enumi})}
\setcounter{enumi}{1}
\tightlist
\item
  Let \(\mathfrak{b}\) be the ideal of \(R\) generated by the \(\beta_{i}\) ; there exists a subset \(J\) of
\end{enumerate}

\[
\mathrm{I} = \{1, 2, \dots , s \}
\]

such that \((\beta_{i})_{i\in\mathbb{J}}\) is a minimal system of generators of b. Then \(J\neq\varnothing\) since \(b\neq0\) . We shall deduce from (9) and (10) that, if \(i\in J\) , \(\alpha_{i}\) is an R-linear combination of the \(\alpha_{j}\) for \(j\neq i\) , which will contradict the minimality of \((\alpha_{1},\ldots,\alpha_{s})\) and will complete the proof.

There exist homogeneous elements \(\gamma_{ji}\) of R ( \(i \in \mathrm{J}, j \in \mathrm{I} - \mathrm{J}\) ) such that

\[
\beta_ {j} = \sum_ {i \in J} \gamma_ {j i} \beta_ {i} (j \in I - J).\tag{11}
\]

Taking (11) into account, formula (10) gives

\[
\sum_ {i \in \mathrm{J}} \beta_ {i} \left(d _ {i k} + \sum_ {j \in \mathrm{I} - \mathrm{J}} \gamma_ {j i} d _ {j k}\right) = 0.\tag{12}
\]

Put

\[
u _ {i k} = d _ {i k} + \sum_ {j \in I - J} \gamma_ {j i} d _ {j k}.\tag{13}
\]

Thus

\[
\sum_ {i \in J} \beta_ {i} u _ {i k} = 0.\tag{14}
\]

Write \(u_{ik} = \sum_{\lambda \in \Lambda} \delta_{ik\lambda} z_\lambda\) , where the \(\delta_{ik\lambda}\) belong to R. Relation (14) implies that \(\sum_{i \in J} \beta_i \delta_{ik\lambda} = 0\) for all \(k\) and \(\lambda\) . If one of the \(\delta_{ik\lambda}\) had a non-zero homogeneous component of degree 0, the preceding equality would imply that one of the \(\beta_i\) ( \(i \in J\) ) is an R-linear combination of the others, contradicting the minimality of \((\beta_i)_{i \in J}\) . Thus \(\delta_{ik\lambda} \in R_+\) and consequently \(u_{ik} \in R_+ S\) for all \(i\) and \(k\) . Thus, there exist \(u_{ikh} \in S\) such that \(u_{ik} = \sum_{h=1}^{s} u_{ikh} \alpha_h\) , in other words, by (13),

\[
d _ {i k} + \sum_ {j \in I - J} \gamma_ {j i} d _ {j k} = \sum_ {h = 1} ^ {s} u _ {i k h} \alpha_ {h}.\tag{$(15)_{ik}$}
\]

Multiply both sides of \((15)_{ik}\) by \(y_{k}\) and add for i in J fixed and \(k = 1, 2, \ldots, r\) ; in view of (9), we find that

\[
\alpha_ {i} + \sum_ {j \in I - J} \gamma_ {j i} \alpha_ {j} = \sum_ {h = 1} ^ {s} \sum_ {k = 1} ^ {r} u _ {i k h} y _ {k} \alpha_ {h}.
\]

Take the homogeneous components of degree \(k_i\) of both sides. Since \(\deg y_k > 0\) , \(\alpha_i\) is an S-linear combination of the \(\alpha_j\) for \(j \neq i\) . Since S is free over R and \(\alpha_1, \ldots, \alpha_s \in \mathrm{R}\) , it follows that \(\alpha_i\) is an R-linear combination of the \(\alpha_j\) with \(j \neq i\) (Commutative Algebra, Chap. I, §3, no. 5, Prop. 9 d)).

COROLLARY. With the assumptions and notation of Th. 3, the product of the characteristic degrees of R is Card(G).

Indeed, \(\operatorname{rk}_{\mathrm{R}}(S) = \operatorname{Card}(G)\) (formula (6), no. 2). The characteristic degrees of S are equal to 1. Then the corollary follows from no. 1, Prop. 2 (iii).

Lemma 2. Let K be a commutative field, V a finite dimensional vector space over K, \(S = \bigoplus_{n \geqslant 0} S_n\) the symmetric algebra of V, s an endomorphism of V, and \(s^{(n)}\) the canonical extension of s to \(S_n\) . Then, with T denoting an indeterminate, we have in K{[}{[}T{]}{]}

\[
\sum_ {n = 0} ^ {\infty} \operatorname{Tr} (s ^ {(n)}) \mathrm{T} ^ {n} = (\det (1 - s \mathrm{T})) ^ {- 1}.
\]

Extending the base field if necessary, we can assume that K is algebraically closed. Let \((e_{1},\ldots,e_{r})\) be a basis of V with respect to which the matrix of s is lower triangular, and let \(\lambda_{1},\ldots,\lambda_{r}\) be the diagonal elements of this matrix. With respect to the basis \((e_{1}^{i(1)}\ldots e_{r}^{i(r)})_{i(1)+\cdots+i(r)=n}\) of \(S_{n}\) , ordered lexicographically, the matrix of \(s^{(n)}\) is lower triangular and its diagonal elements are the products \(\lambda_{1}^{i(1)}\ldots\lambda_{r}^{i(r)}\) . Thus

\[
\operatorname{Tr} \left(s ^ {(n)}\right) = \sum_ {i (1) + \dots + i (r) = n} \lambda_ {1} ^ {i (1)} \dots \lambda_ {r} ^ {i (r)},
\]

and consequently

\[
\begin{array}{l} \sum_ {n = 0} ^ {\infty} \operatorname{Tr} (s ^ {(n)}) \mathrm{T} ^ {n} = \left(\sum_ {n = 0} ^ {\infty} \lambda_ {1} ^ {n} \mathrm{T} ^ {n}\right) \left(\sum_ {n = 0} ^ {\infty} \lambda_ {2} ^ {n} \mathrm{T} ^ {n}\right) \ldots \left(\sum_ {n = 0} ^ {\infty} \lambda_ {r} ^ {n} \mathrm{T} ^ {n}\right) \\ \qquad = (1 - \lambda_ {1} \mathrm{T}) ^ {- 1} (1 - \lambda_ {2} \mathrm{T}) ^ {- 1} \ldots (1 - \lambda_ {r} \mathrm{T}) ^ {- 1} \\ \qquad = (\det (1 - s \mathrm{T})) ^ {- 1}. \end{array}
\]

Lemma 3. Let K, V and S be as in Lemma 2, G a finite group of automorphisms of V, q the order of G, and R the graded subalgebra of S consisting of the elements invariant under G. Assume that K is of characteristic 0. Then the Poincaré series of R is

\[
q ^ {- 1} \sum_ {g \in G} (\det (1 - g T)) ^ {- 1}.
\]

Indeed, the endomorphism \(f = q^{-1} \sum_{g \in G} g^{(n)}\) is a projection of \(S_n\) onto \(R_n\) , so \(\operatorname{Tr}(f) = \dim_K S_n^G\) . Thus, the Poincaré series of \(R\) is

\[
q ^ {- 1} \sum_ {g \in G} \left(\sum_ {n = 0} ^ {\infty} (\operatorname{Tr} g ^ {(n)}) T ^ {n}\right),
\]

and it suffices to apply Lemma 2.

PROPOSITION 3. With the assumptions and notation of Th. 3, let H be the set of pseudo-reflections belonging to G and distinct from 1. Assume that K is of characteristic 0. Then \(\operatorname{Card}(\mathrm{H}) = \sum_{i=1}^{l}(k_i - 1)\) .

By Prop. 3 of the Appendix, we can assume that K is algebraically closed. For any \(g \in G\) , let \(\lambda_{1}(g), \ldots, \lambda_{l}(g)\) be its eigenvalues. Since every \(g \in G\) is diagonalizable (Appendix, Prop. 2), g = 1 if and only if all of the \(\lambda_{i}(g)\) are equal to 1, and \(g \in H\) if and only if the number of \(\lambda_{i}(g)\) equal to 1 is l - 1 (we then denote by \(\lambda(g)\) the eigenvalue distinct from 1). Prop. 1 of no. 1 and Lemma 3 prove that

\[
q \prod_ {i = 1} ^ {l} (1 - T ^ {k _ {i}}) ^ {- 1} = \sum_ {g \in G} (\det (1 - g T)) ^ {- 1}\tag{16}
\]

in \(\mathbf{K}[[\mathbf{T}]]\) , and hence in \(\mathbf{K}(\mathbf{T})\) . Consequently, we have in \(\mathbf{K}(\mathbf{T})\)

\[
q \frac {(1 - T) ^ {l - 1}}{\prod_ {i = 1} ^ {l} (1 - T ^ {k _ {i}})} = \frac {1}{1 - T} + \sum_ {g \in H} \frac {1}{1 - \lambda (g) T} + \sum_ {g \neq 1, g \notin H} \frac {(1 - T) ^ {l - 1}}{\det (1 - g T)},
\]

which can be written

\[
\begin{array}{c} \frac {q - \prod_ {i = 1} ^ {l} (1 + T + \cdots + T ^ {k _ {i} - 1})}{(1 - T) \prod_ {i = 1} ^ {l} (1 + T + \cdots + T ^ {k _ {i} - 1})} \\ = \sum_ {g \in H} \frac {1}{1 - \lambda (g) T} + \sum_ {g \neq 1, g \notin H} \frac {(1 - T) ^ {l - 1}}{\det (1 - g T)}. \end{array}\tag{17}
\]

It follows that \(q - \prod_{i=1}^{l} (1 + T + \cdots + T^{k_i - 1})\) vanishes for T = 1, so \(q = k_1 k_2 \ldots k_l\) , which we knew already by the Cor. of Th. 3. This granted, let \(Q(T)\) be the polynomial \((1 - T)^{-1}(q - \prod_{i=1}^{l} (1 + T + \cdots + T^{k_i - 1}))\) . Differentiating the equality \((1 - T)Q(T) = q - \prod_{i=1}^{l} (1 + T + \cdots + T^{k_i - 1})\) and putting T = 1, we see that \(-Q(1)\) is the value for T = 1 of

\[
\begin{array}{l} - \frac {d}{d t} \left(\prod_ {i = 1} ^ {l} (1 + T + \dots + T ^ {k _ {i} - 1})\right) \\ = - \sum_ {i = 1} ^ {l} \left((1 + 2 T + \dots + (k _ {i} - 1) T ^ {k _ {i} - 2}) \prod_ {j \neq i} (1 + T + \dots + T ^ {k _ {j} - 1})\right) \end{array}
\]

so

\[
Q (1) = \sum_ {i = 1} ^ {l} \frac {\left(k _ {i} - 1\right) k _ {i}}{2} \prod_ {j \neq i} k _ {j} = \left(\prod_ {j = 1} ^ {l} k _ {j}\right) \left(\sum_ {i = 1} ^ {l} \frac {k _ {i} - 1}{2}\right).
\]

Returning to (17), we have on the other hand

\[
Q (1) = \left(\prod_ {j = 1} ^ {l} k _ {j}\right) \left(\sum_ {g \in H} \frac {1}{1 - \lambda (g)}\right).
\]

Thus,

\[
\sum_ {i = 1} ^ {l} \frac {k _ {i} - 1}{2} = \sum_ {g \in H} \frac {1}{1 - \lambda (g)}.\tag{18}
\]

Now the elements of \(\mathbf{G}\) that leave fixed the points of a given hyperplane leave stable a complementary line of the hyperplane (Appendix, Prop. 2), and so form a cyclic subgroup \(\mathbf{G}'\) of \(\mathbf{G}\) by Algebra, Chap. V, §11, no. 1, Th. 1. Let \(t\) be the order of \(\mathbf{G}'\) ; the values of \(\lambda(g)\) for \(g \in \mathbf{G}'\) are \(\theta, \theta^2, \ldots, \theta^{t-1}\) with \(\theta\) a primitive \(t\) th root of unity. We have \(\frac{1}{1-\theta^i} + \frac{1}{1-\theta^{t-i}} = 1\) , so \(\sum_{g \in \mathbf{G}', g \neq 1} \frac{1}{1-\lambda(g)} = \frac{1}{2}(t-1) = \frac{1}{2}\mathrm{Card}(\mathrm{H} \cap \mathrm{G}')\) . The equality (18) thus proves the proposition.

Remark. When \(\mathbf{K} = \mathbf{R}\) , \(\mathbf{G}\) is a Coxeter group and \(\mathbf{H}\) is the set of reflections belonging to \(\mathbf{G}\) , we know (§ 3) that the elements of \(\mathbf{H}\) are in one-to-one correspondence with the walls of \(\mathbf{V}\) .

PROPOSITION 4. With the assumptions and notation of Th. 3, assume that K is of characteristic \(\neq 2\) . Then \(-1 \in G\) if and only if the characteristic degrees \(k_{1}, \ldots, k_{l}\) of R are all even.

Let \(f\) be the automorphism of the algebra \(S\) that extends the automorphism \(-1\) of \(V\) . Then \(f(z) = (-1)^{\deg z}z\) for all homogeneous \(z\) in \(S\) . Thus, if \(-1 \in G\) , every homogeneous element of odd degree of \(R\) is zero, and the \(k_i\) are all even. Conversely, if the \(k_i\) are all even, every element of \(R\) is invariant under \(f\) , and Galois theory shows that \(-1 \in G\) .

